\section{From the original defect to the fifth-line error}
\label{sec:web-primitive-rate}

All norms and integrals in this section are normalized by volume. We
work on the scalar-minus-four Yamabe representative, with $\beta\ge2$,
the original actual conformal monopole supply, $p\ge A>0$, and
$a^2\ge T_+$. No density bound is imposed until the final subsection.
Write
\[
 x=\kappa_\beta+\beta^{-2},\qquad
 b=\fint |W^-|^2,\qquad
 \rho_*=\frac43,\qquad m_*=\frac{128}{81}.
\]
We work with the actual positive spinor $\Phi$, its cubic opposite-spinor
field $\tau$, and the quantities
\[
 \rho=|\Phi|^2,\quad q=|\nabla\Phi|^2,\quad
 \zeta=\langle\Phi,\nabla\Phi\rangle,\quad m=|\tau|^2,
 \quad \nu=2\rho^2q-m\ge0.
\]
The connections in derivatives are the induced actual connections.
The strong-current estimates of Section~\ref{sec:web-strong-current}
provide a fixed $0<\epsilon_0\le1$. Write
\[
 \gamma=\frac{\epsilon_0}{1+\epsilon_0}\le\frac12.
\]
The error $\mathfrak e_D$ below is the full sum of primitive errors in
the fifth-line estimate of Section~\ref{sec:web-density-rate}: it
includes both the $H^1$ and the $L^4$ differences between the bounded
auxiliary field and $\tau$, the amplitude errors, the full prolonged
first-order equation, the Dirac error, the curvature fit, and the
adapted positive-spinor connection errors. In particular, it is not
an independently adjustable error parameter.

\begin{proposition}\label{prop:web-primitive-rate}
For sufficiently small $x$, the auxiliary field can be chosen so that
\[
 \mathfrak e_D\le Cx^{\gamma/16},
\]
where $C$ is independent of $b$. Consequently, on the actual endpoint
class,
\begin{equation}\label{eq:web-polynomial-density}
 \kappa_\beta+\beta^{-2}\ge c(1+b)^{-24/\gamma}.
\end{equation}
\end{proposition}

\begin{proof}
We give all transfers required by the primitive error list. The
previously established exact budgets are
\begin{align*}
 \|\rho-\rho_*\|_2^2+\|\zeta\|_2^2+
 \|\Ric^\circ\|_2^2&\le Cx,\\
 \fint\frac{(S-2m_*)^2}{S}&\le C\sqrt{x},\qquad
 S=m_*+m+\nu=m_*+2\rho^2q,\\
 \|f-K\|_2^2&\le C\sqrt{x},\qquad
 K=-\frac98\frac{\rho^3Q}{S},\qquad |Q|^2=\frac{m^2}{8},\\
 0\le\rho\le2,\qquad \nu&\le4\rho|\zeta|^2.
\end{align*}
Here $f$ is the actual anti-self-dual determinant curvature in the
normalization of the polynomial curvature identity, and $Q$ is the
quadratic Hopf tensor of $\tau$. The strong-current tail estimate is
\[
 \fint_{q>L^2}q^2\le CL^{-2\epsilon_0}+C\sqrt{x}.
\]
Split at $L=x^{-1/(2+2\epsilon_0)}$. The boundedness of $\Phi$ and
$|\zeta|^2\le Cq$ give
\[
 \fint |\zeta|^2q+\fint|\zeta|^4\le Cx^\gamma,
 \qquad \|\nu\|_2\le Cx^{\gamma/2}.
\]
The weighted variance of $S$, together with $\fint S\le C$, implies
$\|m-m_*\|_1\le Cx^{1/4}$. On $m\le M$ use
$(m-m_*)^2\le CM|m-m_*|$; on its complement use the tail estimate.
Taking $M=x^{-1/[4(1+\epsilon_0)]}$ yields
\begin{equation}\label{eq:web-mass-rate}
 \|m-m_*\|_2\le Cx^{\gamma/8}.
\end{equation}
The full cubic Dirac formula, including its axial cancellation, gives
\[
 \|D\tau\|_2^2\le C\left(x+\fint|\zeta|^2q\right),
 \qquad \|D\tau\|_2\le Cx^{\gamma/2}.
\]

For the phase-curvature transfer set
$h(\rho,S)=2-\frac98\rho^6/S$. On $0\le\rho\le2$, $S\ge m_*$,
this function is bounded and Lipschitz, and
$h(\rho_*,2m_*)=0$. The required product estimate follows from
\[
 |hQ|\le C|m-m_*|+Cm_*|h|,
\]
not from an $L^2$ product estimate. Thus
\[
 \|hQ\|_2\le C\bigl(\|m-m_*\|_2+\|\rho-\rho_*\|_2+
                 \|S-2m_*\|_2\bigr)\le Cx^{\gamma/8}.
\]
The polynomial curvature identity implies
$\|\rho^2(d\operatorname{Im}\zeta)^-\|_2\le Cx^{\gamma/8}$.
Fix $0<c_0<\rho_*$. The volume fraction of $\{\rho<c_0\}$ is at
most $Cx$, and splitting $q\le R$ and $q>R$ gives
\[
 \fint_{\rho<c_0}q^2
 \le C\bigl(xR^2+R^{-\epsilon_0}+\sqrt{x}\bigr).
\]
At $R=x^{-1/(2+\epsilon_0)}$ this suffices to remove the weight,
because $\epsilon_0/[2(2+\epsilon_0)]\ge\gamma/8$.
Indeed, $|d\zeta|\le C(q+\rho|F^{S^+}|)$ and $F^{S^+}$ is a
bounded model plus an $L^2$ error. Stokes' identity on the closed
manifold recovers the other chirality. Consequently
\[
 \|d\operatorname{Im}\zeta\|_2+
 \|q-4/9\|_2+\|\Delta\rho\|_2\le Cx^{\gamma/8}.
\]
In the Bochner identities for $\operatorname{Im}\zeta$ and $d\rho$,
the Ricci error satisfies
\[
 \left|\fint(\Ric+g)(\zeta,\zeta)\right|
 \le Cx^{1/2}\|\zeta\|_4^2\le Cx^{(1+\gamma)/2}.
\]
It follows that
\begin{equation}\label{eq:web-full-prolongation-rate}
 \|\nabla\zeta\|_2+\|N_\tau\|_2\le Cx^{\gamma/8}.
\end{equation}
Here $N_\tau$ is the full polynomial prolonged equation, not merely
its Dirac trace.

We next construct a bounded auxiliary section without losing the
rate. Write the actual Lichnerowicz equation as
\[
 D^*D\tau=L\tau-(1+\Gamma)\tau+e\tau,
 \qquad L=\nabla^*\nabla,\qquad
 \Gamma=\frac{27}{32}\frac{\rho^3m}{S},
 \qquad \|e\|_2\le Cx^{1/4}.
\]
In particular $0\le\Gamma\le C$. Fix $m_*<M_1<M_2$, choose a
nondecreasing smooth cutoff $\psi(m)$ equal to zero below $M_1$
and one above $M_2$, and put $\chi=\psi^2$. Then
$m^2|\chi'|^2\le C\chi$. Testing by $\chi\tau$ gives
\begin{align*}
 \fint\chi|\nabla\tau|^2+\tfrac12\fint\chi'|dm|^2
 &=\operatorname{Re}\fint\langle D\tau,D(\chi\tau)\rangle
   +\fint(1+\Gamma)\chi m
   -\operatorname{Re}\fint\chi\langle e\tau,\tau\rangle.
\end{align*}
The second term on the left is nonnegative. The cutoff derivative on
the right obeys
\[
 |D\tau|\,|d\chi|\sqrt m
 \le\varepsilon\chi|\nabla\tau|^2+C_\varepsilon|D\tau|^2.
\]
On $m\ge M_1$, one has $m\le C(m-m_*)^2$. Absorption gives
\begin{equation}\label{eq:web-tail-gradient-rate}
 \fint_{m\ge M_2}|\nabla\tau|^2
 \le C\bigl(\|D\tau\|_2^2+\|m-m_*\|_2^2+
                 \|e\|_2\|m\|_2\bigr)
 \le Cx^{\gamma/4}.
\end{equation}
Choose a fixed smooth radial cap $u=T(\tau)$ whose change is supported
in this high-amplitude region. Then
\[
 |u|\le C,\qquad \|u-\tau\|_{H^1}+\|Du\|_2\le Cx^{\gamma/8}.
\]
Write $D^*D=L+\mathcal V$. Since $u$ is pointwise collinear with
$\tau$, $(\mathcal V+2)u=(1-\Gamma)u+eu$, so
$\|(\mathcal V+2)u\|_2\le Cx^{\gamma/8}$. Set
$v=3(I+L)^{-1}u$. The identity
\[
 (I+L)(v-u)=-D^*Du+(\mathcal V+2)u
\]
and an energy test yield
\begin{equation}\label{eq:web-resolvent-rate}
 \|v-\tau\|_{H^1}+\|(L-2)v\|_2+\|Dv\|_2
 \le Cx^{\gamma/8},\qquad |v|+|Lv|\le C.
\end{equation}
The last bound is the connection maximum principle.

The $L^4$ transfer is essential. Boundedness and the $L^2$ estimate
give $\|v-u\|_4^4\le Cx^{\gamma/4}$. On the cap region,
$m^2\le C(m-m_*)^2$, hence also
$\|u-\tau\|_4^4\le Cx^{\gamma/4}$. Therefore
\begin{equation}\label{eq:web-fourth-norm-transfer}
 \|v-\tau\|_4\le Cx^{\gamma/16}.
\end{equation}
The auxiliary mass $|v|^2$ differs from $m_*$ in $L^2$ by
$O(x^{\gamma/8})$, by comparison with the bounded cap. The full
first-order error transfers by the $H^1$ bound. For the Hopf tensors,
\[
 \|Q_v-Q_\tau\|_2
 \le C(\|v\|_4+\|\tau\|_4)\|v-\tau\|_4
 \le Cx^{\gamma/16}.
\]
The bounded-coefficient decomposition already used for $hQ$ proves
$\|K+(27/32)Q_\tau\|_2\le Cx^{\gamma/8}$, and hence
\[
 \|-iF_A^-+(27/32)Q_v\|_2\le Cx^{\gamma/16}.
\]
On a fixed positive-amplitude region, the positive-spinor phase
coefficient is $\operatorname{Im}\zeta/\rho$, with $L^2$ error
$O(x^{1/2})$. The off-diagonal connection error is bounded by
$C(|v-\tau|+|\zeta|)$, by the actual cubic algebraic formula.
The cutoffs keep all denominators bounded. These estimates cover
every term in $\mathfrak e_D$ and prove its asserted rate.

The fifth-line positive-potential estimate excludes any actual
endpoint sequence for which
$\mathfrak e_D(1+b)^{3/2}\to0$. If
\eqref{eq:web-polynomial-density} failed, a sequence would satisfy
$x(1+b)^{24/\gamma}\to0$, and therefore
\[
 \mathfrak e_D(1+b)^{3/2}
 \le C\bigl(x(1+b)^{24/\gamma}\bigr)^{\gamma/16}\longrightarrow0.
\]
This is the required contradiction.
\end{proof}

\subsection{A direct resolvent alternative}
The cap-gradient estimate can also be bypassed. Let
$(4+L)v=6T(\tau)$ and $w=v-\tau$. Then
\[
 (4+L)w=6(T\tau-\tau)-D^*D\tau-(\Gamma-1)\tau+e\tau.
\]
On $m\le2m_*+1$, the bounded rational coefficient is Lipschitz;
on the complement, $m\le C(m-m_*)^2$. Thus
\[
 \|(\Gamma-1)\tau\|_2
 \le C\bigl(\|\rho-\rho_*\|_2+\|m-m_*\|_2+\|\nu\|_2\bigr)
 \le Cx^{\gamma/8}.
\]
When the equation is tested by $w$, the curvature error requires only
\[
 \left|\fint\langle e\tau,w\rangle\right|
 \le\|e\|_2\|\tau\|_4\|w\|_4\le Cx^{1/4};
\]
indeed $\|w\|_4\le\|v\|_\infty+\|\tau\|_4\le C$.
Absorption proves $\|w\|_{H^1}^2\le Cx^{\gamma/4}$, and the
bounded-cap argument again gives \eqref{eq:web-fourth-norm-transfer}.
No Sobolev constant enters either construction.

\subsection{All parameters at fixed topological density}
For $\chi/A\le B$ and $\kappa_\beta\le1$, the Euler budget gives
$1+b\le C(1+B)$. Put $N=24/\gamma$ and $m_B=c(1+B)^{-N}$.
Equation~\eqref{eq:web-polynomial-density} implies
$\kappa_\beta+\beta^{-2}\ge m_B$. For
$\beta\ge(2/m_B)^{1/2}$ the second term is at most $m_B/2$.
For the remaining parameters the finite-parameter estimate supplies
the matching lower bound after decreasing the universal constant.
The exact conformal transfer of Section~\ref{sec:finite-density-endpoint}
therefore yields
\[
 \fint|\Ric_g+g|^2+\frac{D_\beta(g)}{V_g}
 \ge c'(1+B)^{-N}\frac{(\beta-2)^2}{\beta^2}.
\]
The exponent is not claimed to be optimal. The estimate still permits
$b$ to grow at or faster than $x^{-\gamma/24}$; it is a quantitative
density-dependent endpoint theorem, not a density-independent one.
